%============================================================
%  outlook.tex  –  Chapter 8
%============================================================
\chapter{Outlook}
\label{ch:outlook}

% ---------------------------------------------------------------
% WRITING TIPS – OUTLOOK
% ---------------------------------------------------------------
% Goal: Show you understand the field deeply enough to know what
%       comes next. This chapter demonstrates scientific maturity.
%       It should be forward-looking, speculative (but grounded),
%       and inspiring.
%
% Suggested content (pick the most relevant):
%
%  8.1  Algorithmic Improvements
%       - Implement the full RLE (Reduced Linear Equation) method of
%         Purvis & Bartlett for more stable Hessian inversion.
%       - Trust-region or line-search augmentation to Newton-Raphson
%         to avoid overshooting.
%       - Exploit sparsity of ⟨ij|ab⟩ integrals for large basis sets.
%       - GPU acceleration for the 4-index integral transformation.
%
%  8.2  Method Extensions
%       - Extension to CCSD or CCSD(T) reference (OVOS beyond MP2).
%       - OVOS for excited states: optimise virtual space for EOM-CC.
%       - Restricted (RHF) formalism for closed-shell molecules
%         (simpler, avoids UHF spin contamination).
%       - State-averaged OVOS for multi-reference problems.
%       - Frozen natural orbitals (FNO) vs. OVOS: systematic comparison.
%
%  8.3  Quantum Computing Applications
%       - Feed OVOS-reduced orbital set directly into VQE or QPE.
%       - Benchmark qubit counts and circuit depths on real hardware.
%       - Combine OVOS with active-space selection tools
%         (e.g., AVAS, DMET).
%       - Explore OVOS within quantum embedding theories.
%
%  8.4  Larger and More Complex Systems
%       - Scale to transition-metal complexes, biological molecules.
%       - Periodic systems (solids) using k-point sampling.
%       - Basis set extrapolation combined with orbital truncation.
%
%  8.5  Validation and Benchmarking
%       - Compare OVOS with pair natural orbitals (PNO) and
%         domain-based local pair natural orbitals (DLPNO).
%       - Benchmark against CCSD(T)/CBS as the gold standard.
%       - Test robustness for molecules with strong correlation
%         (e.g., stretched bonds, diradicals).
%
% Tips:
%   - Be specific: "future work should…" is less compelling than
%     "a promising direction is X because Y."
%   - Cite recent papers that your work connects to.
%   - Keep it to 2–3 pages: it is an outlook, not a new proposal.
%   - Do NOT promise results you have not verified – use conditional
%     language ("could", "may", "is expected to").
% ---------------------------------------------------------------

\section{Algorithmic Improvements}
% ---------------------------------------------------------------
% [SECTION INTRO – 8.1]
%   "The current OVOS implementation is a correct but unoptimised
%    proof-of-concept. This section identifies the most impactful
%    algorithmic improvements that would increase numerical stability
%    and computational efficiency."
%
% [WHAT TO COVER]
%   - Full RLE (Reduced Linear Equation) method of Purvis \& Bartlett
%     \cite{Purvis1982}: solves the Newton system iteratively with a
%     Krylov subspace; avoids explicit Hessian inversion.
%   - Trust-region or line-search augmentation: prevent overshooting
%     in regions of large $|\Delta R|$.
%   - Exploit 8-fold symmetry of ERIs: halve memory and integral transform cost.
%   - GPU acceleration: Hessian assembly loops are embarrassingly parallel.
%   - Direct OVOS: compute only the integrals needed at each iteration
%     (avoids storing the full $N_4$ ERI tensor).
%
% [LEAD-IN TO §8.2]
%   "Beyond algorithmic improvements, the method itself can be extended
%    to higher levels of theory, as described in the following section."
% ---------------------------------------------------------------

\section{Method Extensions}
% ---------------------------------------------------------------
% [SECTION INTRO – 8.2]
%   "The OVOS framework developed here is formulated at the MP2 level
%    within an unrestricted reference. This section surveys the most
%    promising extensions to higher correlated levels and alternative
%    electronic structure contexts."
%
% [WHAT TO COVER]
%   - OVOS for CCSD/CCSD(T): minimise the CCSD energy over orbital
%     rotations (\cite{Pitonak2009, Rolik2011} have explored this).
%   - Restricted (RHF) OVOS: cleaner for closed-shell molecules;
%     avoids UHF spin contamination; reduces N by factor of 2.
%   - Excited states: optimise virtual space for EOM-CCSD or TDDFT responses.
%   - Multi-reference: state-averaged OVOS for near-degenerate systems.
%   - Comparison with FNO (frozen natural orbitals) and
%     PNO (pair natural orbitals): systematic benchmark study.
%
% [LEAD-IN TO §8.3]
%   "The most direct application of the results of this thesis is
%    to quantum computing, where OVOS-reduced orbital spaces offer
%    tangible qubit savings, as described in the following section."
% ---------------------------------------------------------------

\section{Quantum Computing Applications}
% ---------------------------------------------------------------
% [SECTION INTRO – 8.3]
%   "The truncation of the virtual orbital space by OVOS is directly
%    relevant to quantum chemistry on near-term and fault-tolerant
%    quantum devices. This section outlines three concrete directions
%    for leveraging OVOS in quantum computing applications."
%
% [LEAD-IN TO §8.3.1] (already has detailed tips below)
%   Start the section with 2-3 sentences contextualising why
%   quantum computing and OVOS are natural partners, then introduce
%   the three subsections.
% ---------------------------------------------------------------
% ---------------------------------------------------------------
% WRITING TIPS – QUANTUM COMPUTING (expanded)
% ---------------------------------------------------------------
%  8.3.1  OVOS Orbitals as VQE Reference State
%       - The thesis demonstrates (Theory §3.10) that OVOS MOs give
%         better overlap with the correlated ground state than HF MOs.
%       - Future work: run actual VQE(UCCSD) calculations in PySCF or
%         Qiskit starting from OVOS vs HF orbitals; compare:
%           * number of UCCSD parameters required for a target accuracy,
%           * circuit depth,
%           * barren-plateau susceptibility.
%       - Benchmark on small molecules (H₂, LiH) accessible to current
%         hardware or high-fidelity classical simulators.
%
%  8.3.2  Brueckner Orbitals and T1 Reduction for NISQ Devices
%       - On NISQ devices, shallow circuits are essential.
%       - If OVOS ≈ Brueckner orbitals, T1 amplitudes are suppressed
%         → UCCSD reduces to UCCD → shallower circuit.
%       - Quantify T1 reduction: compute the Frobenius norm of
%         T1 amplitudes in HF vs OVOS basis for each test molecule.
%       - This provides a concrete metric of how much OVOS helps NISQ.
%
%  8.3.3  Qubit Count Reduction and Fault-Tolerant Perspectives
%       - State the qubit saving per molecule/basis as a table.
%       - For fault-tolerant QPE, fewer logical qubits → fewer T gates
%         → lower error-correction overhead.
%       - Combine OVOS with AVAS or DMET for further active-space
%         reduction.
% ---------------------------------------------------------------
\subsection{OVOS Orbitals as a VQE Reference State}
\label{sec:outlook_vqe}

\subsection{Brueckner Orbitals and Circuit Depth on NISQ Devices}
\label{sec:outlook_brueckner_nisq}

\subsection{Qubit Count Reduction and Fault-Tolerant Perspectives}

\section{Larger Molecular Systems}
% ---------------------------------------------------------------
% [SECTION INTRO – 8.4]
%   "This thesis benchmarks OVOS on small molecules (up to 18 atoms);
%    extending the method to larger and more complex systems would
%    demonstrate its practical utility for real-world applications
%    in catalysis, drug discovery, and materials science."
%
% [WHAT TO COVER]
%   - Transition-metal complexes: strong correlation; test limits of
%     single-reference OVOS.
%   - Biological systems: amino acids, small peptides.
%   - Periodic systems: solid-state MP2 (using k-point sampling);
%     OVOS as a natural partner to plane-wave CC methods.
%   - Basis-set extrapolation: combine OVOS truncation with CBS limit.
%
% [LEAD-IN TO §8.5]
%   "Alongside scaling to larger systems, rigorous benchmarking
%    against state-of-the-art methods is needed to establish
%    the accuracy and reliability of OVOS, as discussed
%    in the final section."
% ---------------------------------------------------------------

\section{Benchmarking and Validation}
% ---------------------------------------------------------------
% [SECTION INTRO – 8.5]
%   "The ultimate test of OVOS is how well it compares to established
%    orbital truncation methods and how reliably it performs across
%    challenging chemical systems. This section outlines the most
%    important benchmarking studies for future work."
%
% [WHAT TO COVER]
%   - Compare OVOS with FNO (frozen natural orbitals): same N',
%     which gives better energy recovery?
%   - Compare with DLPNO-CCSD(T)/CBS: the practical gold standard.
%   - Robustness at stretched geometries (bond-breaking):
%     OVOS at the MP2 level may fail near conical intersections.
%   - Conformational sampling: is OVOS stable across different geometries
%     of the same molecule?
%   - OpenMolcas, PySCF, ORCA: benchmarking against multiple codes
%     removes implementation bias.
%
% [CLOSING SENTENCE FOR ENTIRE OUTLOOK]
%   ``The methods introduced in this thesis represent an entry point
%    into a rich programme of future research that bridges classical
%    orbital optimisation and quantum computational chemistry."
% ---------------------------------------------------------------
